\originalchapter{8.05 2012测试：对称性与三维振子}
原件为Fall 2012 Test, 2012年10月24日12:00--13:30, 五题100分, 闭卷禁计算器; Fall 2013 OCW归档. 以下为\textbf{AI独立解答}.
\originalsection{第1题：十个判断（20分）}
\begin{exercise}
判断: (1)宇称$P\ket x=\ket{-x}$既Hermitian又幺正; (2)$\ket{-x}=-\ket x$; (3)任意$N\times N$矩阵有$N$个独立本征矢; (4)有限维最高能量$E_N$为任意归一态能量均值上界. 随后三项为光滑有限实偶势恰有五束缚态: (5)最高束缚态四节点; (6)非无穷点可同时$\psi=\psi'=0$; (7)激发偶态可在0为0. 最后三项明确为自旋$1/2$, $R=\exp(\ii\pi S_x/\hbar)$: (8)幺正; (9)可对角化; (10)$R=2\ii S_x/\hbar$.
\end{exercise}
\begin{solution}
依次\textbf{T,F,F,T,T,F,F,T,T,T}, 最后三项按卷中自旋$1/2$约定. (1)$P^2=I=P^\dagger P$; (2)两个位置标签不同, 不是同一矢的负号; (3)Jordan块反例; (4)能量均值是本征值的概率加权; (5)一维节点定理; (6)二阶常微分方程零初值唯一为零解; (7)偶态必有$\psi'(0)=0$, 若$\psi(0)=0$也只能零解; (8,9)Hermitian$S_x$的虚指数为正规幺正矩阵; (10)自旋$1/2$时$R=\e^{\ii\pi\sigma_x/2}=\ii\sigma_x$. 对一般自旋$j$最后一个等式不成立, 不能脱离该课程卷的自旋$1/2$范围引用.
\end{solution}
\originalsection{第2题：算符作用的正交分量（15分）}
\begin{exercise}
对归一态与Hermitian$\Omega$, $\Delta\Omega\neq0$, 证明$\Omega\ket\Psi=\langle\Omega\rangle\ket\Psi+\Delta\Omega\ket{\Psi_\perp}$, 其中新矢归一且正交.
\end{exercise}
\begin{solution}
定义$\ket{\Psi_\perp}=(\Omega-\langle\Omega\rangle)\ket\Psi/\Delta\Omega$. 与初态内积为$(\langle\Omega\rangle-\langle\Omega\rangle)/\Delta\Omega=0$; 范数平方为$\langle(\Omega-\langle\Omega\rangle)^2\rangle/(\Delta\Omega)^2=1$. 所以离开原态方向的大小恰是测量宽度. 若宽度0, 态已是本征态, 不需要定义此垂直单位矢.
\end{solution}
\originalsection{第3题：极化恒等式（15分）}
\begin{exercise}
(a)化简$\tfrac14[\bra{u+v}T\ket{u+v}-\bra{u-v}T\ket{u-v}-\ii\bra{u+\ii v}T\ket{u+\ii v}+\ii\bra{u-\ii v}T\ket{u-\ii v}]$. (b)证明复空间中若每个$\bra wT\ket w=0$, 则$T=0$, 并解释实二维90度旋转反例.
\end{exercise}
\begin{solution}
令$A=\bra uT\ket v,B=\bra vT\ket u$, 前两个差为$2(A+B)$, 后两个差带权后为$2(A-B)$, 和除4就是$A$. 若所有对角二次形式都0, 则任意$u,v$有$\bra uT\ket v=0$, 取$u=Tv$给$\norm{Tv}^2=0$, 所以零算符. 实空间没有$\ii v$这两项, 只能得到$A+B$, 反对称90度旋转确可使$v\cdot Tv=0$却非零. 这不要求$T$Hermitian.
\end{solution}
\originalsection{第4题：高次非谐势（25分）}
\begin{exercise}
$V=\hbar^2x^{2n}/(2mL^{2n+2})$,$n\geq2$整数. (a)量纲估基态. (b)用$\psi=\e^{-b^2x^2}$,$b>0$, $c_k=\int\e^{-2y^2}y^{2k}\dd y$, 求变分函数, 不要求优化. (c)画$n\to\infty$势与预期基态能量. (d)给定最优$bL\sim\sqrt n$, 描述试探态是否越好.
\end{exercise}
\begin{solution}
唯一能量尺度$E_L=\hbar^2/(2mL^2)$, 因而$E_g=C_nE_L$. 用$y=bx$归一积分为$c_0/b$, $\langle x^{2n}\rangle=(c_n/c_0)b^{-2n}$. 动能由导数$\psi'=-2b^2x\psi$给$\langle T\rangle=4b^4\hbar^2\langle x^2\rangle/(2m)=\hbar^2b^2/(2m)$, 因$c_1/c_0=1/4$. 所以
\[
F(b)=E_L\left[(bL)^2+\frac{c_n/c_0}{(bL)^{2n}}\right].
\]
极限$|x|<L$势趋0,$|x|>L$趋无限, 端点单个值不影响Dirichlet阱极限; 基态趋宽$2L$盒的$E_L\pi^2/4$, 波函数$\cos(\pi x/(2L))$（内）. Gaussian最优宽度$1/b\sim L/\sqrt n$却缩向0, 比真实分布越来越窄, 因为Gaussian无限尾受高次势强烈惩罚. 试探族越来越差, 不是优化参数保证任何$n$都接近精确态.
\begin{center}\begin{tikzpicture}[x=1.7cm,y=.65cm,>=Stealth]
\draw[->](-1.6,0)--(1.6,0)node[right]{$x/L$};\draw[main](-1,3.4)--(-1,0)--(1,0)--(1,3.4);
\draw[main,domain=-1:1,samples=80]plot(\x,{cos(90*\x)});\node[align=center,font=\footnotesize,text width=3cm]at(0,2.6){实线：盒基态\\点线：Gaussian};
\draw[dotted,domain=-1:1,samples=100]plot(\x,{exp(-16*\x*\x)});\node[below]at(-1,0){$-1$};\node[below]at(1,0){$1$};\end{tikzpicture}\end{center}
\end{solution}
\originalsection{第5题：三维谐振子的角动量（25分）}
\begin{exercise}
三维等频振子有$a_i^\dagger,a_i,N_i$,$i=1,2,3$. (a)用产生算符列基态与前两激发组, 能量与简并. (b)令$a_L=(a_1+\ii a_2)/\sqrt2,a_R=(a_1-\ii a_2)/\sqrt2$, 写$H$, 说明$\{N_L,N_R,N_3\}$完备并求能量. (c)用新产生算符重新列三组. (d)$L_z=\hbar(N_R-N_L)$, 求$[H,L_z]$, 判断$\{H,L_z\}$完备. (e)$L_+=\sqrt2\hbar(a_3^\dagger a_L-a_R^\dagger a_3)$是否Hermitian, 求第二激发组同时被$L_z,L_+$湮灭的唯一态.
\end{exercise}
\begin{solution}
总数$N=0$有$\ket0$, 能量$3\hbar\omega/2$, 一重. $N=1$三态$a_i^\dagger\ket0$, 能量$5\hbar\omega/2$. $N=2$六态$(a_i^\dagger)^2\ket0$及$a_i^\dagger a_j^\dagger\ket0$（$i<j$）, 能量$7\hbar\omega/2$, 未归一允许. 改用$L,R,3$时三组逐项为$\ket0$; $a_L^\dagger\ket0,a_R^\dagger\ket0,a_3^\dagger\ket0$; 第二激发组六态为
\[
\begin{gathered}
(a_L^\dagger)^2\ket0,\quad(a_R^\dagger)^2\ket0,\quad(a_3^\dagger)^2\ket0,\\
a_L^\dagger a_R^\dagger\ket0,\quad a_L^\dagger a_3^\dagger\ket0,\quad a_R^\dagger a_3^\dagger\ket0.
\end{gathered}
\]. 因变换幺正维数不变. $H=\hbar\omega(N_L+N_R+N_3+3/2)$,$E=\hbar\omega(n_L+n_R+n_3+3/2)$, 三独立占有数唯一标记乘积态, 为完备组.

$H,L_z$均为数算符线性组合, 故对易, 却不完备: $a_L^\dagger a_R^\dagger\ket0$和$(a_3^\dagger)^2\ket0$都有$N=2,m=0$. $L_+$伴随为$L_-$, 非Hermitian. 在这二态的组合$[A(a_3^\dagger)^2+B a_L^\dagger a_R^\dagger]\ket0$上作用$L_+$, 得$\sqrt2\hbar(-2A+B)a_R^\dagger a_3^\dagger\ket0$. 所以$B=2A$, 归一态
\[
\ket\psi=\frac{(a_3^\dagger)^2+2a_L^\dagger a_R^\dagger}{\sqrt6}\ket0
=\frac{(a_1^\dagger)^2+(a_2^\dagger)^2+(a_3^\dagger)^2}{\sqrt6}\ket0.
\]
$m=0$且为最高权态只能是$\ell=0$, 故全角动量为0, 不是只有$z$分量为0.
\end{solution}
